\SourceLeaf{359}%
geodätische Linie bekannter Länge und Richtung verbunden ist, ist seit Langem auf vielfältige Weise behandelt worden. Um aus den verschiedenen Methoden eine für den eigenen Gebrauch geeignete zu wählen, muss man einige Umstände berücksichtigen. Es kommt zum Beispiel darauf an, ob man die Aufgabe nur für einen oder einige wenige konkrete Fälle lösen muss oder für sehr viele. Im letzteren Fall ist es wichtig, dass die Methode die abschließende Rechnung jeweils möglichst bequem und übersichtlich macht, auch wenn ihre Anwendung vielleicht erst gewisse allgemeine Vorarbeiten erfordert. Ebenso wichtig ist, ob man alle Ergebnisse einer ausgedehnten trigonometrischen Vermessung in der Form geografischer Länge und Breite und ausschließlich \textit{nur} in \textit{dieser} Form verlangt oder ob die Lage sämtlicher Punkte daneben auch in einer anderen Form angegeben wird, etwa in rechtwinkligen Koordinaten. Im letzteren Fall ist es weniger nötig, die geografische Lage mit allerhöchster Genauigkeit anzugeben.

Die von \textsc{Dusejour}, \textsc{Legendre}, \textsc{Delambre} und anderen angegebenen Formeln berücksichtigen nur die erste Potenz der Abplattung. Das wird in der Praxis kaum ins Gewicht fallen, da die Abplattung des Erdsphäroids nun einmal nur ein kleiner Bruch ist. Die neue Methode ist völlig unabhängig davon, wie klein die Abplattung ist. Das soll hier aber nicht als praktisch wichtiger Vorzug geltend gemacht werden. Die besseren jener Methoden mögen in den meisten Fällen eine ausreichende Genauigkeit bieten, auch wenn ein mathematisch hinreichender Nachweis dafür fehlt. Dagegen lässt sich sagen, dass die neue Methode bei gegebenen Voraussetzungen eine bequemere Rechnung ergibt, deren wesentlicher Inhalt sich auf erheblich kleinerem Raum darstellen lässt. Die von \textsc{Bessel} im Jahr 1825 angegebene Lösung trägt das Gepräge großer mathematischer Vollendung und setzt auch keineswegs voraus, dass der Abstand der beiden Punkte im Vergleich zu den Abmessungen des ganzen Erdsphäroids klein ist. In theoretischer Hinsicht ist das zweifellos ein Vorzug dieser Methode. Bei der Beurteilung ihres praktischen Werts sind aber folgende Umstände zu berücksichtigen. Die Methode unterscheidet überhaupt nicht zwischen großen und kleineren Entfernungen, sondern erfordert in allen Fällen gleich lange Rechnungen,
